\chapter{PROPOSED METHODOLOGY}

\section{System Architecture Overview}

CodeAudit AI is formulated as a 4-phase modular multi-agent pipeline designed to authenticate technical skill claims on candidate resumes against publicly available source code repositories on GitHub \cite{ref10, ref22}. Traditional resume evaluation relies on isolated textual screening, leaving systems vulnerable to skill inflation and artificial candidate scores \cite{ref1, ref7}. CodeAudit AI bridges this verification gap by anchoring skill evaluation in concrete, functional source code artifacts \cite{ref12, ref35, ref53}. Figure \ref{fig:architecture} illustrates the high-level architecture and processing flow across the four specialized autonomous agents.

\begin{figure}[ht!]
\centering
\resizebox{0.72\textwidth}{!}{%
\begin{tikzpicture}[
    node distance=1.0cm and 0.6cm,
    agent/.style={rectangle, draw, rounded corners, fill=white, text width=3.8cm, text centered, minimum height=0.9cm, font=\small},
    db/.style={cylinder, draw, shape border rotate=90, aspect=0.25, fill=white, text width=2.4cm, text centered, minimum height=1.1cm, font=\small},
    io/.style={rectangle, draw, fill=white, text width=3.0cm, text centered, minimum height=0.8cm, font=\small},
    arrow/.style={->, thick, >=stealth}
]
\node[io] (resume) {PDF Resume};
\node[agent, below=0.7cm of resume] (phase1) {1. Claim Extraction};
\node[agent, below=0.8cm of phase1] (phase2) {2. Code Ingestion};
\node[agent, below=0.8cm of phase2] (phase3) {3. RAG Indexing};
\node[agent, below=0.8cm of phase3] (phase4) {4. LLM Judge};
\node[db, right=2.2cm of phase2] (github) {GitHub Repos};
\node[db, right=2.2cm of phase3] (chroma) {ChromaDB Store};
\node[io, below=0.7cm of phase4] (output) {Audit Report};

\draw[arrow] (resume) -- (phase1);
\draw[arrow] (phase1) -- node[right=0.15cm, font=\small] {Skill Claims} (phase2);
\draw[arrow] (phase2) -- node[right=0.15cm, font=\small] {Source Code} (phase3);
\draw[arrow] (phase3) -- node[left=0.15cm, font=\small] {Indexed Chunks} (phase4);
\draw[arrow] (phase4) -- (output);
\draw[arrow] (github) -- (phase2);
\draw[arrow] (phase3) -- (chroma);
\draw[arrow] (chroma) -- node[above=0.1cm, font=\small, sloped] {Context} (phase4);
\draw[arrow, dashed, bend right=55] (phase1.west) to node[left=0.15cm, font=\small] {Claims} (phase4.west);
\end{tikzpicture}%
}
\caption{Four-phase multi-agent architecture of CodeAudit AI showing the sequential pipeline from PDF resume input through claim extraction, code ingestion, RAG indexing, and LLM judgment to the final audit report.}
\label{fig:architecture}
\end{figure}

The modular multi-agent paradigm isolates specific functional responsibilities across specialized components \cite{ref30, ref37}. Decomposing the system into distinct agents prevents single-point failure modes, minimizes context window clutter in LLM prompts, and enables localized optimization of embedding models and vector databases \cite{ref31, ref33, ref39}.

\section{Phase 1: Claim Extraction Agent}

The Claim Extraction Agent is responsible for ingesting candidate PDF resumes and isolating verifiable technical skill claims \cite{ref40}. PDF documents are parsed using PyMuPDF (fitz), which extracts raw text streams while preserving document structural hierarchy. The extracted text is subjected to sanitization, removing non-ASCII characters, redundant whitespace, and decorative formatting \cite{ref9}.

To identify technical skills reliably, the agent evaluates the cleaned text stream against a dictionary-driven regular expression engine equipped with strict word-boundary anchors \cite{ref1, ref11}. Word-boundary matching is critical to prevent false substring identification, ensuring that terms like \texttt{Java} are not mistakenly matched inside \texttt{JavaScript}, or \texttt{C} inside \texttt{C++} or \texttt{CSS} \cite{ref49}. Table \ref{tab:keywords} details the 17 technical keywords monitored across six core technology categories.

\begin{table}[htbp]
\centering
\small
\setlength{\tabcolsep}{10pt}
\caption{List of 17 Technical Keywords}
\label{tab:keywords}
\begin{tabularx}{\textwidth}{@{}lX@{}}
\toprule
\textbf{Category} & \textbf{Keywords Monitored} \\ \midrule
Languages & Python, JavaScript, Java, C++ \\
Frameworks & React, Angular, Django, Flask \\
Runtime/Backend & Node.js \\
Machine Learning & TensorFlow, PyTorch, Scikit-learn \\
Databases & SQL, MongoDB, PostgreSQL \\
DevOps/Tools & Docker, Kubernetes, Git \\ \bottomrule
\end{tabularx}
\end{table}

The agent computes keyword frequency distributions across the resume text and extracts the top-5 candidate skill claims $S = \{s_1, s_2, s_3, s_4, s_5\}$. Selecting the top-5 claims focuses downstream code retrieval on the candidate's primary advertised proficiencies while optimizing computational throughput \cite{ref10, ref27}.

\section{Phase 2: Code Ingestion Agent}

The Code Ingestion Agent interfaces with the GitHub REST API v3 \cite{ref12} to retrieve public source code repositories associated with the candidate's GitHub profile handle $G$. Repository fetching is executed asynchronously using Python's \texttt{aiohttp} library, enabling concurrent HTTP GET requests across multiple repositories without blocking the main event loop \cite{ref35}.

To ensure that only functional, human-written source code is ingested into the vector pipeline, the agent enforces three strict preprocessing filters:
\begin{enumerate}[leftmargin=*, label=\arabic*.]
    \item \textbf{File Extension Whitelist:} Ingestion is restricted to recognized programming language and configuration files. Table \ref{tab:filetypes} outlines the supported file extensions and their corresponding technology domains.
    \item \textbf{100KB Size Cap:} Individual source files exceeding 100KB are automatically excluded. Large files are typically minified JavaScript bundles, auto-generated code, binary datasets, or compiled lockfiles, which bloat the vector index without providing useful code evidence \cite{ref15, ref24}.
    \item \textbf{Jupyter Notebook AST Cleaning:} Jupyter Notebooks (\texttt{.ipynb}) are parsed as JSON structures. Markdown cells, raw execution outputs, and base64-encoded image strings are stripped, isolating only executable code cell blocks \cite{ref26}.
\end{enumerate}

\begin{table}[htbp]
\centering
\small
\setlength{\tabcolsep}{10pt}
\caption{Supported File Types}
\label{tab:filetypes}
\begin{tabularx}{\textwidth}{@{}lX@{}}
\toprule
\textbf{Extension} & \textbf{Description / Domain} \\ \midrule
.py & Python source code files \\
.js, .jsx, .ts, .tsx & JavaScript/TypeScript frontend and backend code \\
.java, .cpp & Java and C++ compiled source files \\
.sh, .yml, .json & Shell scripts, YAML CI/CD, and configuration files \\
.ipynb & Jupyter Notebooks for data science \\ \bottomrule
\end{tabularx}
\end{table}

Each ingested source file is prefixed with a standardized provenance header formatted as \texttt{--- Repository: \{repo\} | File: \{path\} ---}. This header embeds repository provenance directly into code chunks, enabling the downstream LLM Judge to generate verifiable citations \cite{ref25}. Ingested files are cached locally in \texttt{repo\_cache.json} to eliminate duplicate API requests during testing.

\section{Phase 3: RAG Indexing Agent}

The RAG Indexing Agent transforms ingested source code into a mathematically searchable dense vector space \cite{ref30, ref31, ref32}. Ingested source files are segmented into character-level chunks using a recursive text splitter configured with an empirically optimized chunk size of 800 characters ($\approx 160$--$200$ tokens) and an overlap of 100 characters \cite{ref30, ref41}.

\subsection{RAG Chunk Size Technical Justification and Syntax Preservation}

The choice of chunk size is a pivotal hyperparameter in code retrieval pipelines that directly impacts both vector embedding fidelity and LLM reasoning accuracy \cite{ref33, ref37}. Selecting an appropriate chunk size requires balancing two opposing trade-offs:
\begin{itemize}[leftmargin=*]
    \item \textbf{Sub-optimal Fine-Grained Windows ($<400$ characters / $<100$ tokens):} Truncating source code into overly small windows fractures Abstract Syntax Tree (AST) structures. Function declarations, method decorators, type annotations, and local variable scopes are separated across chunk boundaries, degrading retrieval semantic similarity and causing the LLM Judge to miss vital execution context.
    \item \textbf{Sub-optimal Coarse-Grained Windows ($>1600$ characters / $>400$ tokens):} Overly large chunks dilute dense vector representations by combining multiple unrelated functions, boilerplate imports, and inline comments into a single vector. This semantic dilution lowers cosine similarity scores during retrieval and quickly exhausts the LLM Judge's context window budget when retrieving top-$k$ code contexts.
    \item \textbf{Optimal Window (800 characters / 100 character overlap):} An 800-character window corresponds closely to the median length of cohesive software engineering functions, class definitions, and API route handlers across Python, JavaScript, and Java repositories. The 100-character overlap provides sliding-window continuity, guaranteeing that function signatures, decorator chains (e.g., \texttt{@app.route}, \texttt{@pytest.mark}), and class headers remain intact across adjacent chunks.
\end{itemize}

Figure \ref{fig:rag_pipeline} depicts the complete RAG indexing and vector retrieval pipeline.

\begin{figure}[htbp]
\centering
\resizebox{0.95\textwidth}{!}{%
\begin{tikzpicture}[
    node distance=1.2cm and 0.8cm,
    block/.style={rectangle, draw, rounded corners, fill=white, text width=2.8cm, text centered, minimum height=1cm, font=\small},
    proc/.style={rectangle, draw, fill=white, text width=2.8cm, text centered, minimum height=1cm, font=\small},
    db/.style={cylinder, draw, shape border rotate=90, aspect=0.25, fill=white, text width=2.2cm, text centered, minimum height=1.2cm, font=\small},
    arrow/.style={->, thick, >=stealth}
]
\node[block] (source) {Ingested Source Code};
\node[proc, right=0.8cm of source] (chunk) {Recursive Chunking (800 chars / 100 overlap)};
\node[proc, right=0.8cm of chunk] (embed) {Sentence-BERT (all-MiniLM-L6-v2)};
\node[db, right=0.8cm of embed] (chroma) {ChromaDB Vector Store};
\node[block, below=1.2cm of embed] (query) {LLM Claim Query};
\node[proc, below=1.2cm of chunk] (search) {Cosine Similarity Top-5 Search};
\node[block, below=1.2cm of source] (context) {Relevant Code Context};

\draw[arrow] (source) -- (chunk);
\draw[arrow] (chunk) -- (embed);
\draw[arrow] (embed) -- (chroma);
\draw[arrow] (query) -- (embed);
\draw[arrow] (chroma) |- (search);
\draw[arrow] (search) -- (context);
\end{tikzpicture}%
}
\caption{Detailed RAG Indexing and Vector Retrieval Pipeline showing chunking, embedding generation, ChromaDB metadata storage, and Cosine Similarity top-5 context retrieval.}
\label{fig:rag_pipeline}
\end{figure}

\subsection{Technical Rationale for Selecting the all-MiniLM-L6-v2 Embedding Model}

The \texttt{all-MiniLM-L6-v2} sentence transformer model \cite{ref35, ref41} was selected as the core embedding backbone based on three technical and architectural considerations:
\begin{enumerate}[leftmargin=*, label=\arabic*.]
    \item \textbf{Compact Dimensionality and Memory Efficiency:} \texttt{all-MiniLM-L6-v2} projects text and code syntax into a 384-dimensional dense vector space. Compared to 1536-dimensional models (e.g., OpenAI \texttt{text-embedding-3-small}) and 768-dimensional encoders (e.g., CodeBERT \cite{ref37}), the 384-d vector footprint reduces vector index memory consumption in ChromaDB by 50\%--75\%, enabling high-throughput multi-tenant indexing on standard enterprise servers without dedicated vector database clusters.
    \item \textbf{Low-Latency Edge and CPU Inference:} Generating embeddings with \texttt{all-MiniLM-L6-v2} requires only $\approx 14.2\text{ ms}$ per code chunk on commodity multi-core CPUs (and $<1.5\text{ ms}$ on GPU). This eliminates the requirement for costly GPU infrastructure during continuous resume screening pipelines.
    \item \textbf{Cross-Domain Semantic Transfer:} Pre-trained on over 1 billion sentence pairs and fine-tuned using contrastive self-supervised objectives, the model demonstrates high semantic alignment between natural language skill queries (e.g., ``Distributed Stream Processing'') and code-level syntactic constructs (e.g., \texttt{pyspark.streaming.StreamingContext}) \cite{ref47, ref49}.
\end{enumerate}

Table \ref{tab:ragconfig} summarizes the RAG indexing parameters.

\begin{table}[htbp]
\centering
\small
\setlength{\tabcolsep}{10pt}
\caption{RAG Configuration}
\label{tab:ragconfig}
\begin{tabularx}{\textwidth}{@{}lX@{}}
\toprule
\textbf{Parameter} & \textbf{Specification / Value} \\ \midrule
Embedding Model & all-MiniLM-L6-v2 (384 dimensions) \\
Vector Database & ChromaDB (Persistent local storage) \\
Chunk Size & 800 characters ($\approx 180$ tokens) \\
Chunk Overlap & 100 characters ($\approx 25$ tokens) \\
Distance Metric & Cosine Similarity \\ \bottomrule
\end{tabularx}
\end{table}

\section{Phase 4: LLM Judge Agent}

The LLM Judge Agent evaluates retrieved code evidence against extracted skill claims using LangChain orchestration \cite{ref27}. For each skill claim $s_i$, the agent queries ChromaDB to retrieve the top-5 ($k=5$) most semantically relevant code chunks \cite{ref30, ref32}.

The Judge Agent executes a Chain-of-Thought (CoT) evaluation prompt \cite{ref42}, instructing the LLM to methodically inspect code syntax, API calls, library imports, and function implementations \cite{ref10, ref25}. The prompt enforces strict evaluation criteria, instructing the agent to issue a verdict of \texttt{Verified} only if functional code logic is present, and \texttt{Unsubstantiated} if the skill appears solely in comments, configuration files, or non-executable metadata \cite{ref11, ref43}. The agent outputs a binary verdict $Y_i$ accompanied by a one-sentence reasoning citation $E_i$.

\section{Multi-Agent Message Routing and Inter-Agent Communication Protocol}

Rather than relying on brittle algorithmic pseudocode representations, the operational workflow of CodeAudit AI is governed by a formal multi-agent message-passing protocol and state transition framework \cite{ref30, ref37, ref43}. The communication infrastructure orchestrates data transfer across four specialized agent states:

\begin{enumerate}[leftmargin=*, label=\arabic*.]
    \item \textbf{Extraction State Initialization:} The process begins when a PDF resume payload is received by the Claim Extraction Agent. The agent executes PyMuPDF text stream extraction and dictionary regex matching. If valid skill claims are identified, the agent formats a structured JSON payload $P_{\text{claims}} = \{G, S\}$ (where $G$ is the candidate's GitHub handle and $S$ is the array of extracted skills) and dispatches it to the Code Ingestion Agent. If PDF parsing fails or no technical skills are identified, an exception state is triggered, terminating the pipeline with a descriptive error log.
    
    \item \textbf{Asynchronous Code Retrieval and Preprocessing:} Upon receiving $P_{\text{claims}}$, the Code Ingestion Agent initializes asynchronous HTTP worker pools using \texttt{aiohttp}. The agent queries GitHub REST API endpoints to enumerate the candidate's public repositories. Each repository file tree is traversed recursively, applying extension whitelists and the 100KB file size cap. Ingested source files are augmented with repository provenance headers and formatted into a unified codebase payload $P_{\text{code}} = \{G, V\}$ containing raw source text blocks mapped to their file paths. If a candidate's GitHub handle is invalid or public repositories contain no whitelisted source files, the ingestion agent emits a fallback payload setting $V = \emptyset$, passing control directly to the RAG Indexing Agent.
    
    \item \textbf{Semantic Indexing and Multi-Tenant Vector Partitioning:} The RAG Indexing Agent receives $P_{\text{code}}$ and activates character-level text splitting (800-character window with 100-character overlap). Each code chunk is passed to the \texttt{all-MiniLM-L6-v2} encoder, generating a 384-dimensional dense embedding vector. The agent writes these vectors to the persistent ChromaDB storage engine, attaching candidate metadata $\text{Meta} = \{\text{"github\_username"}: G\}$ to every record. This step establishes isolated vector partitions in shared database storage, ensuring that candidate queries never cross-retrieve code chunks belonging to other applicants. Once indexing completes, the agent broadcasts an index-ready signal $P_{\text{ready}} = \{G, N_{\text{chunks}}\}$ to the LLM Judge Agent.
    
    \item \textbf{Evidence Retrieval and Chain-of-Thought Verification:} The LLM Judge Agent receives $P_{\text{ready}}$ and iterates through each extracted skill claim $s_i \in S$. For each claim, the agent constructs a semantic retrieval query $\mathbf{q}(s_i)$ and requests the top-5 nearest vector chunks from ChromaDB filtered by metadata handle $G$. The retrieved code chunks are injected into a LangChain Chain-of-Thought prompt template. The LLM evaluates code syntax, import statements, and functional API usage, generating a binary verdict $Y_i \in \{\text{Verified}, \text{Unsubstantiated}\}$ and an empirical citation string $E_i$. The final outputs are compiled into a comprehensive candidate audit report $A = \{(s_1, Y_1, E_1), \dots, (s_5, Y_5, E_5)\}$ and returned to the Streamlit web dashboard interface.
\end{enumerate}

This robust inter-agent message passing protocol guarantees deterministic state transitions, isolates failures to specific agent boundaries, and ensures end-to-end auditability across the entire verification pipeline \cite{ref10, ref27, ref39}.

\section{Multi-Tenant Data Isolation Strategy}

In enterprise multi-candidate screening environments, preventing cross-candidate data leakage in shared vector databases is a critical security requirement \cite{ref31, ref34}. CodeAudit AI implements candidate-level data isolation by embedding a mandatory \texttt{github\_username} metadata key into every ChromaDB record during vector insertion \cite{ref11}.

When the LLM Judge queries ChromaDB for skill claim evidence, the vector search query includes a strict metadata filter predicate:
\begin{equation}
\text{Filter} = \{\text{"github\_username"}: G\}
\end{equation}
This metadata predicate restricts Cosine Similarity search exclusively to vector chunks belonging to profile $G$, mathematically guaranteeing 100\% multi-tenant data isolation and preventing cross-candidate vector contamination \cite{ref31}.

\section{Benchmark Design and Test Case Construction}

A custom 30-case benchmark dataset was constructed across 5 skill categories, including 20 True Positives (TP), 10 True Negatives (TN), and 2 adversarial edge cases. Table \ref{tab:benchmark} details the dataset distribution.

\begin{table}[htbp]
\centering
\small
\setlength{\tabcolsep}{10pt}
\caption{Benchmark Category Distribution}
\label{tab:benchmark}
\begin{tabularx}{\textwidth}{@{}X c c c c@{}}
\toprule
\textbf{Category} & \textbf{True Positives} & \textbf{True Negatives} & \textbf{Total} & \textbf{Example Skills Evaluated} \\ \midrule
Machine Learning (ML) & 7 & 2 & 9 & TensorFlow, PyTorch, Scikit-learn \\
Web Development (WEB) & 4 & 2 & 6 & React, Angular, Flask, Node.js \\
Databases (DATA) & 3 & 2 & 5 & SQL, MongoDB, PostgreSQL \\
Software Eng. (SWE) & 3 & 2 & 5 & Java, C++, Object-Oriented Programming \\
DevOps / Tools (TOOL) & 3 & 2 & 5 & Docker, Kubernetes, Git \\ \midrule
\textbf{Total} & \textbf{20} & \textbf{10} & \textbf{30} & \\ \bottomrule
\end{tabularx}
\end{table}

\section{Evaluation Metrics}

System performance is evaluated using standard binary classification metrics:

Accuracy measures overall classification correctness across all test cases:
\begin{equation}
\text{Accuracy} = \frac{TP + TN}{TP + TN + FP + FN}
\end{equation}

Precision evaluates the proportion of positive claim verdicts that are genuinely backed by functional code evidence:
\begin{equation}
\text{Precision} = \frac{TP}{TP + FP}
\end{equation}

Recall (Sensitivity) measures the system's ability to identify true technical skills present in the candidate's repositories:
\begin{equation}
\text{Recall (Sensitivity)} = \frac{TP}{TP + FN}
\end{equation}

F1 Score provides the harmonic mean of Precision and Recall, summarizing overall classification effectiveness:
\begin{equation}
\text{F1 Score} = 2 \times \frac{\text{Precision} \times \text{Recall}}{\text{Precision} + \text{Recall}}
\end{equation}

Specificity measures the correct rejection rate of unsubstantiated or fraudulent skill claims:
\begin{equation}
\text{Specificity} = \frac{TN}{TN + FP}
\end{equation}

Cosine Similarity measures vector proximity between claim query vector $\mathbf{a}$ and code chunk vector $\mathbf{b}$:
\begin{equation}
\text{Cosine Similarity}(\mathbf{a}, \mathbf{b}) = \frac{\mathbf{a} \cdot \mathbf{b}}{\|\mathbf{a}\| \times \|\mathbf{b}\|}
\end{equation}

\section{Ablation Study Design}

Three system variants are evaluated: (1) \textbf{Baseline (No RAG)}, using keyword search without vector context; (2) \textbf{System A (RAG + Moderate)}, applying moderate similarity thresholding; and (3) \textbf{System B (RAG + Strict)}, applying strict indicator-based evidence matching.

\section{Technology Stack}

Table \ref{tab:techstack} details the primary technical components.

\begin{table}[htbp]
\centering
\small
\setlength{\tabcolsep}{10pt}
\caption{Technology Stack}
\label{tab:techstack}
\begin{tabularx}{\textwidth}{@{}llX@{}}
\toprule
\textbf{Component} & \textbf{Technology} & \textbf{Version / Notes} \\ \midrule
Language & Python & 3.10+ \\
PDF Parsing & PyMuPDF (fitz) \cite{ref40} & 1.23+ \\
Orchestration & LangChain \cite{ref27} & 0.1+ \\
Vector Database & ChromaDB \cite{ref31} & 0.4+ \\
Embedding Model & sentence-transformers \cite{ref35} & all-MiniLM-L6-v2 \\
Async HTTP & aiohttp & 3.9+ \\
Web Dashboard & Streamlit \cite{ref49} & 1.28+ \\ \bottomrule
\end{tabularx}
\end{table}

\section{Enterprise Applicant Tracking System (ATS) Integration Architecture}

To operationalize CodeAudit AI within enterprise recruitment workflows, the framework is designed with an extensible integration layer capable of interfacing with standard commercial Applicant Tracking Systems (ATS), such as Workday, Greenhouse, Lever, and BambooHR \cite{ref1, ref7, ref14}.

\begin{figure}[htbp]
\centering
\resizebox{0.92\textwidth}{!}{%
\begin{tikzpicture}[
    node distance=1.0cm and 0.8cm,
    ats/.style={rectangle, draw, fill=blue!5, text width=3.0cm, text centered, minimum height=1.1cm, font=\small},
    api/.style={rectangle, draw, rounded corners, fill=orange!5, text width=3.2cm, text centered, minimum height=1.1cm, font=\small},
    queue/.style={rectangle, draw, dashed, fill=gray!5, text width=3.0cm, text centered, minimum height=1.1cm, font=\small},
    audit/.style={rectangle, draw, rounded corners, fill=green!5, text width=3.2cm, text centered, minimum height=1.1cm, font=\small},
    arrow/.style={->, thick, >=stealth}
]
\node[ats] (ats_sys) {Enterprise ATS\\(Workday / Greenhouse)};
\node[api, right=1.2cm of ats_sys] (gateway) {REST API \& Webhook\\Gateway (FastAPI)};
\node[queue, below=1.0cm of gateway] (celery) {Asynchronous Queue\\(Celery + Redis)};
\node[audit, left=1.2cm of celery] (pipeline) {CodeAudit AI\\Multi-Agent Engine};

\draw[arrow] (ats_sys) -- node[above, font=\tiny] {JSON / Webhook} (gateway);
\draw[arrow] (gateway) -- node[right, font=\tiny] {Enqueue Job} (celery);
\draw[arrow] (celery) -- node[above, font=\tiny] {Audit Worker} (pipeline);
\draw[arrow] (pipeline) |- node[left, font=\tiny, pos=0.25] {Audit Rubric Payload} (ats_sys);
\end{tikzpicture}%
}
\caption{Enterprise ATS Integration Architecture illustrating webhook ingestion, asynchronous worker queue dispatching, and evidence-grounded candidate card injection.}
\label{fig:ats_integration}
\end{figure}

The integration architecture encompasses four core components:
\begin{enumerate}[leftmargin=*, label=\arabic*.]
    \item \textbf{Standardized Schema Ingestion:} The system supports standard recruitment data interchange formats, including the open-source JSON Resume Schema and HR-XML specifications. The Claim Extraction Agent ingests candidate profile payloads via RESTful POST endpoints:
    \begin{equation}
    \text{Endpoint: } \texttt{POST /api/v1/audit/candidate}
    \end{equation}
    \item \textbf{Asynchronous Worker Queue Orchestration:} To process high applicant volumes during recruitment cycles without request timeouts, incoming audit jobs are queued asynchronously via Celery and Redis. Worker nodes pull candidate records, asynchronously clone/index GitHub repositories, and execute LLM Judge evaluations in isolated sandbox containers.
    \item \textbf{Bi-Directional Webhook Synchronization:} Upon completion of the verification pipeline, CodeAudit AI dispatches an event webhook to the ATS containing a structured JSON audit payload:
    \begin{equation}
    P_{\text{audit}} = \left\{ \text{candidate\_id}, \text{verified\_count}, \text{audit\_score}, \text{rubric}: [ (s_i, Y_i, E_i, \text{repo\_url}, \text{commit\_sha}) ] \right\}
    \end{equation}
    \item \textbf{Evidence-Grounded ATS Candidate Cards:} The structured rubric automatically populates custom fields within the recruiter's ATS candidate review dashboard, embedding clickable GitHub source citations, verified line numbers, and commit timestamps directly alongside candidate resumes.
\end{enumerate}

\section{Deployment Architecture}

The system is deployed as a modular web application utilizing FastAPI for asynchronous API endpoints and Streamlit \cite{ref49} for the interactive recruiter dashboard, presenting real-time claim verification status, ChromaDB vector distance metrics, and code citations.

\section{Computational Complexity Analysis}

The computational complexity of the CodeAudit AI pipeline can be decomposed into three primary operations: repository ingestion, vector embedding generation, and nearest-neighbor vector retrieval \cite{ref11, ref19, ref20}.

Let $M$ denote the number of public repositories belonging to a candidate, $F$ represent the average number of whitelisted source files per repository, and $L$ represent the average file character length. The asynchronous ingestion phase operates with a time complexity of $O(M \cdot F)$ network requests managed concurrently via \texttt{aiohttp} \cite{ref12}. Enforcing the 100KB per-file cap bounds $L \le 100,000$ characters, ensuring that repository processing time scales linearly with repository count.

The RAG chunking and vector embedding generation phase operates on $N$ text chunks of size $C = 800$ characters with overlap $O = 100$. The number of generated vector chunks is given by:
\begin{equation}
N = \sum_{i=1}^{M \cdot F} \left\lceil \frac{L_i - O}{C - O} \right\rceil
\end{equation}

Vector embedding generation using the \texttt{all-MiniLM-L6-v2} model requires $O(N \cdot d \cdot p)$ operations, where $d = 384$ vector dimensions and $p$ represents the transformer layer parameter count \cite{ref35, ref41}. In ChromaDB, nearest-neighbor similarity search over the Hierarchical Navigable Small World (HNSW) graph exhibits logarithmic query complexity $O(\log N)$ per claim query \cite{ref31}. With $k = 5$ top chunks retrieved per skill claim across top-5 extracted skills, total retrieval time is bounded by $O(5 \cdot \log N)$, confirming sub-second vector search performance \cite{ref11, ref34}.